
\documentclass{article}
\usepackage{colortbl}
\usepackage{makecell}
\usepackage{multirow}
\usepackage{supertabular}

\begin{document}

\newcounter{utterance}

\centering \large Interaction Transcript for game `bbh', experiment `disambiguation\_qa', episode 17 with reflexion{-}Qwen3.5{-}9B.
\vspace{24pt}

{ \footnotesize  \setcounter{utterance}{1}
\setlength{\tabcolsep}{0pt}
\begin{supertabular}{c@{$\;$}|p{.15\linewidth}@{}p{.15\linewidth}p{.15\linewidth}p{.15\linewidth}p{.15\linewidth}p{.15\linewidth}}
    \# & \multicolumn{2}{c}{Player} && \multicolumn{2}{c}{Game Master} \\
    \hline

    \theutterance \stepcounter{utterance}  
    & & & \multicolumn{4}{p{0.6\linewidth}}{
        \cellcolor[rgb]{0.9,0.9,0.9}{
            \makecell[{{p{\linewidth}}}]{
                \texttt{\tiny{[P1$\langle$GM]}}
                \texttt{You have attempted this same task before and did not succeed. These are the plans you wrote after those attempts.} \\
\texttt{Trial \#1: New plan: I incorrectly assumed that "his" must refer to the employee because feedback is usually given *about* someone else, overlooking the grammatical rule that a pronoun often refers to the closest preceding male noun (the supervisor). In this specific context, "his stellar performance" could logically describe the supervisor's own excellent work, which the employee is being informed about. Next time, I will not rely solely on pragmatic context to resolve ambiguity; instead, I will strictly evaluate if the sentence structure allows for the pronoun to refer to the subject (the supervisor). If both the supervisor and the employee are plausible antecedents without creating a logical contradiction, I will select option (C) Ambiguous rather than forcing a specific antecedent.} \\
\\ 
\texttt{Trial \#2: New plan: I failed because I over{-}analyzed the pragmatic context of "feedback" and ignored the strict grammatical proximity rule where a possessive pronoun like "his" often defaults to the closest preceding male noun ("the employee"). In the previous attempt, I correctly identified the ambiguity but still output the full analysis text instead of the required concise format "A: (C)". Next time, I will immediately identify the closest male noun as the primary antecedent; if the sentence structure allows the pronoun to refer to that noun without contradiction, I will select that specific option (B) and output only the required format "A: (B)" without any explanatory text.} \\
\\ 
\texttt{Trial \#3: New plan: I failed because I incorrectly assumed the sentence was ambiguous by over{-}relying on the pragmatic context that feedback is usually given *to* someone about *their* performance, ignoring the grammatical rule that a possessive pronoun like "his" typically refers to the closest preceding male noun ("the employee"). In this specific context, it is perfectly logical for the supervisor to be giving feedback on the employee's stellar performance. Next time, I will strictly apply the proximity rule: identify the closest male noun ("the employee") as the antecedent. If the sentence structure allows the pronoun to refer to that noun without contradiction, I will select option (B) and output only the required format "A: (B)" without any explanatory text.} \\
\\ 
\\ 
\texttt{Clarify the meaning of sentences with ambiguous pronouns.} \\
\\ 
\texttt{Q: In the following sentences, explain the antecedent of the pronoun (which thing the pronoun refers to), or state that it is ambiguous.} \\
\texttt{Sentence: The chief told the counselor that they took the day off.} \\
\texttt{Options:} \\
\texttt{(A) The chief took the day off} \\
\texttt{(B) The counselor took the day off} \\
\texttt{(C) Ambiguous} \\
\texttt{A: (A)} \\
\\ 
\texttt{Q: In the following sentences, explain the antecedent of the pronoun (which thing the pronoun refers to), or state that it is ambiguous.} \\
\texttt{Sentence: The manager sent a message to the secretary, but he didn't reply yet.} \\
\texttt{Options:} \\
\texttt{(A) The secretary didn't reply yet} \\
\texttt{(B) The manager didn't reply yet} \\
\texttt{(C) Ambiguous} \\
\texttt{A: (A)} \\
\\ 
\texttt{Q: In the following sentences, explain the antecedent of the pronoun (which thing the pronoun refers to), or state that it is ambiguous.} \\
\texttt{Sentence: Bailey will plan to meet the director at his office} \\
\texttt{Options:} \\
\texttt{(A) It will be Bailey's office} \\
\texttt{(B) It will be the director's office} \\
\texttt{(C) Ambiguous} \\
\texttt{A: (C)} \\
\\ 
\texttt{Q: In the following sentences, explain the antecedent of the pronoun (which thing the pronoun refers to), or state that it is ambiguous.} \\
\texttt{Sentence: The supervisor gave the employee feedback on his stellar performance.} \\
\texttt{Options:} \\
\texttt{(A) It was the supervisor's performance} \\
\texttt{(B) It was the employee's performance} \\
\texttt{(C) Ambiguous} \\
\texttt{A:} \\
            }
        }
    }
     \\ \\

    \theutterance \stepcounter{utterance}  
    & \multicolumn{4}{p{0.6\linewidth}}{
        \cellcolor[rgb]{0.561,0.737,0.561}{
            \makecell[{{p{\linewidth}}}]{
                \texttt{\tiny{[P1$\rangle$GM]}}
                \texttt{A: (B)} \\
            }
        }
    }
    & & \\ \\

    \theutterance \stepcounter{utterance}  
    & & & \multicolumn{2}{p{0.3\linewidth}}{
        \cellcolor[rgb]{0.9,0.9,0.9}{
            \makecell[{{p{\linewidth}}}]{
                \texttt{\tiny{[GM$|$GM]}}
                \texttt{Target: (B)} \\
            }
        }
    }
    & & \\ \\

    \theutterance \stepcounter{utterance}  
    & & & \multicolumn{2}{p{0.3\linewidth}}{
        \cellcolor[rgb]{0.9,0.9,0.9}{
            \makecell[{{p{\linewidth}}}]{
                \texttt{\tiny{[GM$|$GM]}}
                \texttt{game\_result = LOSE} \\
            }
        }
    }
    & & \\ \\

    \theutterance \stepcounter{utterance}  
    & & \multicolumn{4}{p{0.6\linewidth}}{
        \cellcolor[rgb]{0.95,0.95,0.95}{
            \makecell[{{p{\linewidth}}}]{
                \texttt{\tiny{[P1$\leftrightarrow$P]}}
                \texttt{Post{-}episode Reflexion (not sent to GM):} \\
\texttt{New plan: I failed because I incorrectly assumed the sentence was ambiguous by over{-}relying on the pragmatic context that feedback is usually given *to* someone about *their* performance, ignoring the grammatical rule that a possessive pronoun like "his" typically refers to the closest preceding male noun ("the employee"). In this specific context, it is perfectly logical for the supervisor to be giving feedback on the employee's stellar performance. Next time, I will strictly apply the proximity rule: identify the closest male noun ("the employee") as the antecedent. If the sentence structure allows the pronoun to refer to that noun without contradiction, I will select option (B) and output only the required format "A: (B)" without any explanatory text.} \\
            }
        }
    }
    & \\ \\

\end{supertabular}
}

\end{document}
